\begin{helper}[Induction state for Algorithm (B)]
Consider the decreasing \noderef{ExtensionBaseAlgorithmBTrace} on the
positions below \(b\), with position \(b-1-a\) at step \(a\).  Assume that
\(\varrho_1\) is a proper partial edge coloring leaving exactly \(e_0\)
uncolored and that, for every trace step \(a\),
\(\mathtt{trace.gamma}(a)=\gamma(\mathtt{trace.position}(a))\) and
\(\mathtt{trace.delta}(a)=\delta(\mathtt{trace.position}(a))\).  Also assume
the Case 1 color chain, distinctness, initial missing-color, minimality,
initial strong-interchangeability, exact stored-pair, and terminal
stored-union hypotheses.  Then every coloring in the trace is proper and
leaves exactly \(e_0\) uncolored.  Immediately before every step, its delta
is missing and its gamma is present at \(u\), both colors are present at the
terminal support \(v_n\), and its pair is strongly interchangeable.  After
that step, every unprocessed slot in the exact Case 1 index set remains
strongly interchangeable under the new coloring.  The analogous assertion
holds at the final coloring for any slot absent from the whole trace.
\end{helper}
\begin{proof}
Write \(r_a=b-1-a\), \(t_a=\mathtt{indices}(r_a)\), and
\(U=\bigcup_{t\in L_1^*}S_t\), where
\(L_1^*=\noderef{ExtensionBaseCaseOneIndices}\).  Define
\(c_0=\gamma_0\) and \(c_{q+1}=\delta_q\).  The chain identity gives
\(\gamma_q=c_q\), and injectivity of
\(\operatorname{Fin.cons}(\gamma_0,\delta)\) makes all \(c_q\) distinct.
The exact trace-field clause of
\noderef{ExtensionBaseAlgorithmBTrace}, together with the assumed trace
gamma and delta equalities, identifies the colors recorded at position
\(r\) with \(\gamma_r\) and \(\delta_r\).  Consequently the pair switched
at position \(r\) is
\(\{c_r,c_{r+1}\}\), and every lower gamma \(c_q\), \(q<r\), differs from
both switched colors.  The stored-pair equalities put both colors of every
processed slot in \(U\).

We induct on the trace step.  By the initial-coloring clause of
\noderef{ExtensionBaseAlgorithmBTrace}, the initial trace coloring is
\(\varrho_1\); hence properness and the unique uncolored edge are the
corresponding hypotheses.  At the first position \(r=b-1\), the chain
identity gives \(\delta_r=\gamma_b\), which is missing at \(u\); minimality
says that every \(\gamma_q\), \(q<b\), is present there by
\noderef{MissingColors} and \noderef{PresentColors}.  Initial strong
interchangeability is supplied at every position below \(b\).  Finally,
the empty intersection at \(v_n\) says that no color of \(U\) is missing
there, so all colors of \(U\), and in particular the current pair, are
present at \(v_n\).

Assume these assertions immediately before the step at \(r\).  Apply
\noderef{ExtensionBaseAlgorithmBSingleStep} to the current proper coloring,
current strong-interchangeability certificate, current missing delta and
present gamma at \(u\), and presence of both colors at \(v_n\).  It supplies
the rooted nontrivial path and its tree-intersection properties.  The two
colors are distinct by the injectivity argument above, so
\noderef{TwoColorComponentSwapProper} makes the component-swapped coloring
proper.  The transition clause of \noderef{ExtensionBaseAlgorithmBTrace}
identifies this coloring with the successor coloring in the trace.
Inspection of \noderef{TwoColorComponentSwap} shows pointwise that an edge
is uncolored after the swap exactly when it was uncolored before: on the
component the two nonempty color values are exchanged, and off it every
value is fixed.  Thus \noderef{IsOneEdgeUncolored} is preserved.

At \(u\), the missing delta means that there is no incident delta edge.
Properness makes the incident gamma edge unique: two distinct such edges
would be adjacent through \(u\) and have the same color, contrary to
\noderef{IsProperEdgeColoring}, using \noderef{AdjacentEdges} and
\noderef{Incident}.  By \noderef{TwoColorIncidenceGraph}, that gamma
edge-node is adjacent to the root vertex-node \(u\), so
\noderef{TwoColorComponent} places it in the rooted component and the swap
changes it to delta.  No incident edge can acquire gamma: this would have
to come either from an old delta edge in the component or from an old gamma
edge outside it, and both possibilities have just been excluded.  Hence
gamma becomes missing and delta becomes present.  Since every lower
\(\gamma_q\) differs from the two switched colors, its incident witness is
fixed.  If a next position exists, its delta is the present step's gamma by
the chain identity, proving the missing-color descent.

At \(v_n\), if its vertex-node is outside the rooted component, every
incident edge is fixed.  Indeed, \noderef{TwoColorIncidenceGraph} and
\noderef{TwoColorComponent} say that an edge-node in the rooted component
has a switched color and, if incident with \(v_n\), is adjacent to the
vertex-node \(v_n\).  Such an incident edge-node would therefore put
\(v_n\) in the component, a contradiction.  Thus no incident edge lies in
the swapped component, and \noderef{TwoColorComponentSwap} fixes all of
them.  If \(v_n\) is inside, the one-step lemma gives distinct incident
gamma and delta edges; the swap exchanges their colors, while every other
color is fixed.  Thus every color of \(U\) remains present at \(v_n\).

It remains to propagate the unprocessed pairs.  Exact enumeration by the
trace, as stated in \noderef{ExtensionBaseAlgorithmBTrace}, writes any
unprocessed slot as \(t_c\).  Since it was not processed at or before step
\(a\), one has \(c>a\), and the strict-descent clause of the same node gives
\(t_c<t_a\).  The induction hypothesis supplies strong
interchangeability for both its pair and the current pair.  Instantiate
\noderef{StrongInterchangeabilityKempePreserved} with their common support
\(v_n\), the strict slot inequality, the current rooted nontrivial path,
and the proper/exactly-one-uncolored assertions.  Its same-support branch
shows that the earlier unprocessed pair remains strongly interchangeable
after the swap.  The exact-field clause of
\noderef{ExtensionBaseAlgorithmBTrace} and the assumed gamma and delta
trace equalities identify any colors named in the conclusion with those at
position \(r_c\).

This completes the induction and proves every asserted pre-step and
post-step clause.  At the last coloring the same survivor argument applies.
In fact the exact-image clause of \noderef{ExtensionBaseAlgorithmBTrace}
says that the trace enumerates every member of \(L_1^*\), so a member absent
from all steps is impossible; the stated final implication is therefore
vacuous, including when \(b=0\), where \(L_1^*=\varnothing\).
\end{proof}
